\section{Related work\label{sec:related}}
%We review related work in terms of treatment effect estimation and semi-supervised learning.
\subsection{Treatment effect estimation}
Treatment effect estimation has been one of the major interests in causal inference and widely studied in various domains.
Matching~\cite{rubin1973matching,abadie2006large} is one of the most basic and commonly used treatment effect estimation techniques. 
It estimates the counterfactual outcomes using its nearby instances, whose idea is similar to that of graph-based semi-supervised learning. 
Both methods assume that similar instances in terms of covariates have similar outcomes.   
% To mitigate the curse of dimensionality problem in matching, the propensity matching method relying on one-dimensional propensity score was proposed~\cite{rosenbaum1983central}. 
% The propensity score is the probability of an instance to get treatment, and is also used for debiasing the treatment biases~\cite{}.
To mitigate the curse of dimensionality and selection bias in matching, the propensity matching method relying on the one-dimensional propensity score was proposed~\cite{rosenbaum1983central,rosenbaum1985constructing}. 
The propensity score is the probability of an instance to get a  treatment, which is modeled using probabilistic models like logistic regression, and has been successfully applied in various domains to estimate treatment effects unbiasedly~\cite{lunceford2004stratification}. 
% Tree-based methods can build an expressive model based on many weak learners~\cite{chipman2010bart,wager2018estimation}.
Tree-based methods such as regression trees and random forests have also been well studied for this problem~\cite{chipman2010bart,wager2018estimation}.
%and they make many subsets of instances and learn propensity tree. 
%They use regression trees divide the covariate space into 
%The leaves of trained trees are uzed to estimate treatment effect~\cite{wager2018estimation}. 
One of the advantages of such models is that they can build quite expressive and flexible models to estimate treatment effects.
% Recently, deep learning-based methods have been successfully applied for treating the effect estimation~\cite{shalit2017estimating,johansson2016learning}. 
Recently, deep learning-based methods have been successfully applied to treatment effect estimation~\cite{shalit2017estimating,johansson2016learning}. 
Balancing neural networks~(BNNs)~\cite{johansson2016learning} aim to obtain balanced representations of a treatment groups and a control group by minimizing the discrepancy between them, such as the Wasserstein distance~\cite{shalit2017estimating}. 
Most recently, some studies have addressed causal inference problems on network-structured data~\cite{guo2020learning,alvari2019less,veitch2019using}.  
Alvari et al. applied the idea of manifold regularization using users activities as causality-based features to detect harmful users in social media~\cite{alvari2019less}. 
Guo et al. considered treatment effect estimation on social networks using graph convolutional balancing neural networks~\cite{guo2020learning}. 
In contrast with their work assuming the network structures are readily available, we do not assume them and considers matching network defined using covariates.

% knn~\cite{lunceford2004stratification}

\subsection{Semi-supervised learning}
% Semi-supervised learning have been a very popular method. Thus it is impossible to cover all works, we explain 
Semi-supervised learning, which exploits both labeled and unlabeled data, is one of the most popular approaches, especially in scenarios when only limited labeled data can be accessed~\cite{bengio2007greedy,hinton2006fast}.
Semi-supervised learning has many variants, and because it is almost impossible to refer to all of them, we mainly review the graph-based regularization methods, known as label propagation or manifold regularization~\cite{zhu2002learning,belkin2006manifold,weston2012deep}. 
Utilizing a given graph or a graph constructed based on instance proximity, graph-based regularization encourages the neighbor instances to have similar labels or outcomes~\cite{zhu2002learning,belkin2006manifold}. 
%This technique is called {\it label propagation} because the key idea behind this technique is that the neighbors of each instance are regularized to have the same label. 
Such idea is also applied to representation learning in deep neural networks~\cite{weston2012deep,yang2016revisiting,kipf2016semi,iscen2019label,bui2018neural}; they encourage nearby instances not only to have similar outcomes, but also have similar intermediate representations, which results in remarkable improvements from ordinary methods. 
One of the major drawbacks of semi-supervised approaches is that label noises in training data can be quite harmful; therefore, a number of studies managed to mitigate the performance degradation~\cite{vahdat2017toward,pal2018label,du2018robust,liu2012robust}.

One of the most related work to our present study is graph-based semi-supervised prediction under sampling biases of labeled data~\cite{zhou2019graph}.
%The graph-based semi-supervised problem under the sampling biases of labeled data is one of the most related works~\cite{zhou2019graph}.
The important difference between this work and ours is that they do not consider intervention and we do not consider the sampling biases of labeled data. 





